\documentclass[10pt,a4paper]{article}
\usepackage[utf8]{inputenc}
\usepackage[T1]{fontenc}
\usepackage[portuguese]{babel}
\usepackage[margin=1.9cm]{geometry}
\usepackage{xcolor}
\usepackage{array}
\setlength{\parindent}{0pt}
\setlength{\parskip}{7pt}
\definecolor{epazul}{RGB}{36,117,154}
\begin{document}
% Gerado de gerar_relatorio.py; editar a fonte comum.
\section*{EP1 - Escalonador de Processos}
\subsection*{Relatório de avaliação do sistema}
Sistemas Operacionais - 2026. Implementação em Rust, edição 2021.

\subsection*{Integrantes}
Cristian de Lima Meira - USP: 16860754

Fernando Silva - USP: 15636702

Julio Domingos - USP: 16836502

\subsection*{Objetivo e resultado}
Avaliar o escalonador de uma máquina fictícia com um processador, registradores X/Y e instruções de atribuição, COM, E/S e SAIDA. A avaliação usa os dez programas oficiais e compara as médias exigidas pelo EP em quinze valores de quantum.

Adotamos quantum 7: é o menor valor testado que alcança pelo menos 85\% da média máxima de instruções por despacho desta carga. Ele atinge 86,84\% dessa média e reduz os despachos em 72,66\% em relação a q1. O limiar de 85\% é uma escolha prática dos autores, não uma prova de ótimo universal.

Os resultados vêm dos eventos dos logs, conferidos contra as instruções e prioridades. PDF e TEX são derivados do mesmo conteúdo em gerar\_relatorio.py. As diferenças de diagramação entre os dois formatos não alteram textos, tabelas ou séries numéricas.

\clearpage
\section*{1. Método e carga de trabalho}
A série principal contém q1 a q10, com passo unitário: dez valores uniformemente distribuídos no intervalo útil de 1 a 10 instruções. Esse intervalo cobre a queda mais acentuada dos despachos e o critério de escolha. Os pontos q12, q15, q20, q25 e q30 são uma exploração complementar, com espaçamento não uniforme.

O limite de um programa é 21 instruções. Os pontos q25 e q30 excedem esse limite e servem para confirmar a saturação; não são necessários para representar novos efeitos de preempção. Prioridades e instruções foram mantidas em todos os casos.

\begin{center}
\small
\begin{tabular}{ccccc}
\hline
\textbf{Arquivo} & \textbf{Programa} & \textbf{Instruções} & \textbf{E/S} & \textbf{Prioridade} \\
\hline
01.txt & TESTE-1 & 10 & 2 & 1 \\
02.txt & TESTE-2 & 21 & 3 & 3 \\
03.txt & TESTE-3 & 14 & 0 & 2 \\
04.txt & TESTE-4 & 12 & 5 & 10 \\
05.txt & TESTE-5 & 11 & 5 & 4 \\
06.txt & TESTE-6 & 15 & 3 & 6 \\
07.txt & TESTE-7 & 14 & 2 & 12 \\
08.txt & TESTE-8 & 21 & 0 & 8 \\
09.txt & TESTE-9 & 15 & 2 & 8 \\
10.txt & TESTE-10 & 6 & 1 & 4 \\
\hline
\end{tabular}
\end{center}
\subsection*{Contagens e validação}
A carga possui 10 processos, 139 instruções e 33 trechos terminados por E/S ou SAIDA. A média natural desses trechos é 139/33 = 4,2121 instruções; o maior tem 21 instruções.

Para cada quantum, o gerador exige igualdade byte a byte dos logs da raiz e de experimentos/. Em seguida reconstitui carga, fila de prontos, créditos, esperas, PC, X/Y, instruções por surto, eventos e resumo. Todos os quinze casos são validados antes da escrita de qualquer saída. Isso confere os artefatos existentes; não substitui uma nova execução do binário Rust.

\clearpage
\section*{2. Resultados e definição das métricas}
P é o número inicial de processos; D é o número de despachos concluídos; I é o total de instruções executadas. A média de trocas do EP é D/P, incluindo a última execução e despachos consecutivos do mesmo PID. Ela não representa apenas mudanças físicas entre processos distintos. E/S e SAIDA contam em I; ociosidade não conta em D ou I.

A média de instruções por quantum é I/D. Definimos também R = (I/D)/q: a fração média do limite de instruções utilizada em cada despacho. R é uma razão de contagens lógicas, não uma medida de utilização da CPU em tempo real.

\begin{center}
\small
\begin{tabular}{cccccc}
\hline
\textbf{q} & \textbf{D} & \textbf{I} & \textbf{D/P} & \textbf{I/D} & \textbf{R (\%)} \\
\hline
1 & 139 & 139 & 13,90 & 1,0000 & 100,00 \\
2 & 79 & 139 & 7,90 & 1,7595 & 87,97 \\
3 & 59 & 139 & 5,90 & 2,3559 & 78,53 \\
4 & 52 & 139 & 5,20 & 2,6731 & 66,83 \\
5 & 44 & 139 & 4,40 & 3,1591 & 63,18 \\
6 & 41 & 139 & 4,10 & 3,3902 & 56,50 \\
7 & 38 & 139 & 3,80 & 3,6579 & 52,26 \\
8 & 38 & 139 & 3,80 & 3,6579 & 45,72 \\
9 & 37 & 139 & 3,70 & 3,7568 & 41,74 \\
10 & 37 & 139 & 3,70 & 3,7568 & 37,57 \\
12 & 35 & 139 & 3,50 & 3,9714 & 33,10 \\
15 & 34 & 139 & 3,40 & 4,0882 & 27,25 \\
20 & 34 & 139 & 3,40 & 4,0882 & 20,44 \\
25 & 33 & 139 & 3,30 & 4,2121 & 16,85 \\
30 & 33 & 139 & 3,30 & 4,2121 & 14,04 \\
\hline
\end{tabular}
\end{center}
Tabela 1. Medições extraídas e validadas a partir dos logs. Nas quinze execuções, P = 10 e I = 139. Médias são calculadas em ponto flutuante; o arredondamento acima é apenas de apresentação.

\clearpage
\section*{3. Efeito do quantum nas médias}
\begin{center}
\setlength{\unitlength}{1pt}
\begin{picture}(470,200)
{\color{gray!25}\put(44.000,32.000){\line(1,0){410}}}
\put(38.000,32.000){\makebox(0,0)[r]{\scriptsize 0}}
{\color{gray!25}\put(44.000,82.000){\line(1,0){410}}}
\put(38.000,82.000){\makebox(0,0)[r]{\scriptsize 5}}
{\color{gray!25}\put(44.000,132.000){\line(1,0){410}}}
\put(38.000,132.000){\makebox(0,0)[r]{\scriptsize 10}}
{\color{gray!25}\put(44.000,182.000){\line(1,0){410}}}
\put(38.000,182.000){\makebox(0,0)[r]{\scriptsize 15}}
\put(44.000,32.000){\line(0,1){150}}
\put(44.000,32.000){\line(1,0){410}}
\put(44.000,19.000){\makebox(0,0){\scriptsize 1}}
\put(100.552,19.000){\makebox(0,0){\scriptsize 5}}
\put(171.241,19.000){\makebox(0,0){\scriptsize 10}}
\put(241.931,19.000){\makebox(0,0){\scriptsize 15}}
\put(312.621,19.000){\makebox(0,0){\scriptsize 20}}
\put(383.310,19.000){\makebox(0,0){\scriptsize 25}}
\put(454.000,19.000){\makebox(0,0){\scriptsize 30}}
\put(44.000,191.000){\makebox(0,0)[l]{\small D/P}}
\put(249.000,3.000){\makebox(0,0){\small Quantum q (instruções)}}
\color{epazul}\linethickness{1pt}
\qbezier(44.000,171.000)(51.069,141.000)(58.138,111.000)
\qbezier(58.138,111.000)(65.207,101.000)(72.276,91.000)
\qbezier(72.276,91.000)(79.345,87.500)(86.414,84.000)
\qbezier(86.414,84.000)(93.483,80.000)(100.552,76.000)
\qbezier(100.552,76.000)(107.621,74.500)(114.690,73.000)
\qbezier(114.690,73.000)(121.759,71.500)(128.828,70.000)
\qbezier(128.828,70.000)(135.897,70.000)(142.966,70.000)
\qbezier(142.966,70.000)(150.034,69.500)(157.103,69.000)
\qbezier(157.103,69.000)(164.172,69.000)(171.241,69.000)
\qbezier(171.241,69.000)(185.379,68.000)(199.517,67.000)
\qbezier(199.517,67.000)(220.724,66.500)(241.931,66.000)
\qbezier(241.931,66.000)(277.276,66.000)(312.621,66.000)
\qbezier(312.621,66.000)(347.966,65.500)(383.310,65.000)
\qbezier(383.310,65.000)(418.655,65.000)(454.000,65.000)
\put(44.000,171.000){\circle*{4}}
\put(58.138,111.000){\circle*{4}}
\put(72.276,91.000){\circle*{4}}
\put(86.414,84.000){\circle*{4}}
\put(100.552,76.000){\circle*{4}}
\put(114.690,73.000){\circle*{4}}
\put(128.828,70.000){\circle*{4}}
\put(142.966,70.000){\circle*{4}}
\put(157.103,69.000){\circle*{4}}
\put(171.241,69.000){\circle*{4}}
\put(199.517,67.000){\circle*{4}}
\put(241.931,66.000){\circle*{4}}
\put(312.621,66.000){\circle*{4}}
\put(383.310,65.000){\circle*{4}}
\put(454.000,65.000){\circle*{4}}
\end{picture}
\end{center}
{\small Figura 1. Média de trocas do EP por processo em função de q.}

A maior redução ocorre nos primeiros valores de quantum. Depois, aumentá-lo reduz cada vez menos os despachos, porque os trechos de CPU terminam em E/S ou SAIDA.

\begin{center}
\setlength{\unitlength}{1pt}
\begin{picture}(470,200)
{\color{gray!25}\put(44.000,32.000){\line(1,0){410}}}
\put(38.000,32.000){\makebox(0,0)[r]{\scriptsize 0}}
{\color{gray!25}\put(44.000,62.000){\line(1,0){410}}}
\put(38.000,62.000){\makebox(0,0)[r]{\scriptsize 1}}
{\color{gray!25}\put(44.000,92.000){\line(1,0){410}}}
\put(38.000,92.000){\makebox(0,0)[r]{\scriptsize 2}}
{\color{gray!25}\put(44.000,122.000){\line(1,0){410}}}
\put(38.000,122.000){\makebox(0,0)[r]{\scriptsize 3}}
{\color{gray!25}\put(44.000,152.000){\line(1,0){410}}}
\put(38.000,152.000){\makebox(0,0)[r]{\scriptsize 4}}
{\color{gray!25}\put(44.000,182.000){\line(1,0){410}}}
\put(38.000,182.000){\makebox(0,0)[r]{\scriptsize 5}}
\put(44.000,32.000){\line(0,1){150}}
\put(44.000,32.000){\line(1,0){410}}
\put(44.000,19.000){\makebox(0,0){\scriptsize 1}}
\put(100.552,19.000){\makebox(0,0){\scriptsize 5}}
\put(171.241,19.000){\makebox(0,0){\scriptsize 10}}
\put(241.931,19.000){\makebox(0,0){\scriptsize 15}}
\put(312.621,19.000){\makebox(0,0){\scriptsize 20}}
\put(383.310,19.000){\makebox(0,0){\scriptsize 25}}
\put(454.000,19.000){\makebox(0,0){\scriptsize 30}}
\put(44.000,191.000){\makebox(0,0)[l]{\small I/D}}
\put(249.000,3.000){\makebox(0,0){\small Quantum q (instruções)}}
\color{epazul}\linethickness{1pt}
\qbezier(44.000,62.000)(51.069,73.392)(58.138,84.785)
\qbezier(58.138,84.785)(65.207,93.731)(72.276,102.678)
\qbezier(72.276,102.678)(79.345,107.435)(86.414,112.192)
\qbezier(86.414,112.192)(93.483,119.483)(100.552,126.773)
\qbezier(100.552,126.773)(107.621,130.240)(114.690,133.707)
\qbezier(114.690,133.707)(121.759,137.722)(128.828,141.737)
\qbezier(128.828,141.737)(135.897,141.737)(142.966,141.737)
\qbezier(142.966,141.737)(150.034,143.220)(157.103,144.703)
\qbezier(157.103,144.703)(164.172,144.703)(171.241,144.703)
\qbezier(171.241,144.703)(185.379,147.923)(199.517,151.143)
\qbezier(199.517,151.143)(220.724,152.895)(241.931,154.647)
\qbezier(241.931,154.647)(277.276,154.647)(312.621,154.647)
\qbezier(312.621,154.647)(347.966,156.505)(383.310,158.364)
\qbezier(383.310,158.364)(418.655,158.364)(454.000,158.364)
\put(44.000,62.000){\circle*{4}}
\put(58.138,84.785){\circle*{4}}
\put(72.276,102.678){\circle*{4}}
\put(86.414,112.192){\circle*{4}}
\put(100.552,126.773){\circle*{4}}
\put(114.690,133.707){\circle*{4}}
\put(128.828,141.737){\circle*{4}}
\put(142.966,141.737){\circle*{4}}
\put(157.103,144.703){\circle*{4}}
\put(171.241,144.703){\circle*{4}}
\put(199.517,151.143){\circle*{4}}
\put(241.931,154.647){\circle*{4}}
\put(312.621,154.647){\circle*{4}}
\put(383.310,158.364){\circle*{4}}
\put(454.000,158.364){\circle*{4}}
\end{picture}
\end{center}
{\small Figura 2. Média de instruções executadas por despacho em função de q.}

Uma E/S ou SAIDA encerra o surto mesmo quando ainda resta quantum. Por isso I/D não acompanha q indefinidamente: o limite desta carga é 4,2121.

\clearpage
\section*{4. Razão utilizada e saturação}
\begin{center}
\setlength{\unitlength}{1pt}
\begin{picture}(470,200)
{\color{gray!25}\put(44.000,32.000){\line(1,0){410}}}
\put(38.000,32.000){\makebox(0,0)[r]{\scriptsize 0,00}}
{\color{gray!25}\put(44.000,67.714){\line(1,0){410}}}
\put(38.000,67.714){\makebox(0,0)[r]{\scriptsize 0,25}}
{\color{gray!25}\put(44.000,103.429){\line(1,0){410}}}
\put(38.000,103.429){\makebox(0,0)[r]{\scriptsize 0,50}}
{\color{gray!25}\put(44.000,139.143){\line(1,0){410}}}
\put(38.000,139.143){\makebox(0,0)[r]{\scriptsize 0,75}}
{\color{gray!25}\put(44.000,174.857){\line(1,0){410}}}
\put(38.000,174.857){\makebox(0,0)[r]{\scriptsize 1,00}}
\put(44.000,32.000){\line(0,1){150}}
\put(44.000,32.000){\line(1,0){410}}
\put(44.000,19.000){\makebox(0,0){\scriptsize 1}}
\put(100.552,19.000){\makebox(0,0){\scriptsize 5}}
\put(171.241,19.000){\makebox(0,0){\scriptsize 10}}
\put(241.931,19.000){\makebox(0,0){\scriptsize 15}}
\put(312.621,19.000){\makebox(0,0){\scriptsize 20}}
\put(383.310,19.000){\makebox(0,0){\scriptsize 25}}
\put(454.000,19.000){\makebox(0,0){\scriptsize 30}}
\put(44.000,191.000){\makebox(0,0)[l]{\small R = (I/D)/q}}
\put(249.000,3.000){\makebox(0,0){\small Quantum q (instruções)}}
\color{epazul}\linethickness{1pt}
\qbezier(44.000,174.857)(51.069,166.268)(58.138,157.678)
\qbezier(58.138,157.678)(65.207,150.933)(72.276,144.187)
\qbezier(72.276,144.187)(79.345,135.827)(86.414,127.467)
\qbezier(86.414,127.467)(93.483,124.863)(100.552,122.260)
\qbezier(100.552,122.260)(107.621,117.490)(114.690,112.720)
\qbezier(114.690,112.720)(121.759,109.686)(128.828,106.651)
\qbezier(128.828,106.651)(135.897,101.985)(142.966,97.320)
\qbezier(142.966,97.320)(150.034,94.475)(157.103,91.631)
\qbezier(157.103,91.631)(164.172,88.650)(171.241,85.668)
\qbezier(171.241,85.668)(185.379,82.473)(199.517,79.279)
\qbezier(199.517,79.279)(220.724,75.107)(241.931,70.936)
\qbezier(241.931,70.936)(277.276,66.069)(312.621,61.202)
\qbezier(312.621,61.202)(347.966,58.635)(383.310,56.069)
\qbezier(383.310,56.069)(418.655,54.063)(454.000,52.058)
\put(44.000,174.857){\circle*{4}}
\put(58.138,157.678){\circle*{4}}
\put(72.276,144.187){\circle*{4}}
\put(86.414,127.467){\circle*{4}}
\put(100.552,122.260){\circle*{4}}
\put(114.690,112.720){\circle*{4}}
\put(128.828,106.651){\circle*{4}}
\put(142.966,97.320){\circle*{4}}
\put(157.103,91.631){\circle*{4}}
\put(171.241,85.668){\circle*{4}}
\put(199.517,79.279){\circle*{4}}
\put(241.931,70.936){\circle*{4}}
\put(312.621,61.202){\circle*{4}}
\put(383.310,56.069){\circle*{4}}
\put(454.000,52.058){\circle*{4}}
\end{picture}
\end{center}
{\small Figura 3. Fração média do quantum utilizada em cada despacho.}

R diminui com q porque os surtos naturais são curtos em boa parte da carga. Esse gráfico compara a média executada com o limite permitido, como pede o enunciado. Uma razão menor não significa, por si só, CPU desperdiçada: o processo libera o processador ao bloquear ou terminar.

D/P e I/D não são critérios independentes: seu produto é I/P = 139/10 = 13,9. Um gráfico entre essas duas médias apenas repete essa identidade; não demonstra um compromisso com responsividade ou tempo de espera.

Sem preempção interna aos trechos, cada um exige um despacho: D mínimo = 33 e I/D máximo = 4,2121. A saturação ocorre em q = 21, o tamanho do maior trecho. Esse valor é deduzido das instruções, não uma execução adicional incluída na tabela. q25 e q30 confirmam o mesmo mínimo observado. Quantum grande mantém a política de créditos; não transforma a fila de prontos em FIFO.

\clearpage
\section*{5. Escolha do quantum}
O critério escolhido é atingir pelo menos 85\% da média máxima de instruções por despacho com o menor quantum testado. Ele combina uma meta explícita de redução dos despachos com um limite menor de instruções consecutivas por execução; não otimiza uma função de tempo real.

\begin{center}
\small
\begin{tabular}{cccc}
\hline
\textbf{q} & \textbf{D} & \textbf{I/D} & \textbf{Fração do máximo (\%)} \\
\hline
6 & 41 & 3,3902 & 80,49 \\
7 & 38 & 3,6579 & 86,84 \\
8 & 38 & 3,6579 & 86,84 \\
\hline
\end{tabular}
\end{center}
q6 alcança 80,49\% do máximo e não atinge a meta. q7 alcança 86,84\% e é o primeiro valor testado que a satisfaz. q8 tem os mesmos D e I/D de q7, mas permite uma instrução adicional consecutiva; portanto q7 é preferível segundo o critério adotado.

\subsection*{Benefícios e custo da escolha}
Em q7, D/P = 3,80 e I/D = 3,6579. Há 72,66\% menos despachos que em q1 e 35,59\% menos que em q3. Em contrapartida, q7 usa 15,15\% mais despachos que o mínimo observado de 33 (3,30 por processo).

O tamanho médio natural dos trechos é 4,2121, não 7. A escolha de q7 decorre da meta de 85\%, e não de uma suposta igualdade com esse tamanho médio. q3 é o valor da entrada original fornecida; os resultados não permitem inferir que seus autores priorizaram responsividade.

\subsection*{Limites da conclusão}
Não foram medidos tempo de resposta, espera, throughput em tempo real ou custo de salvar/carregar contexto. Quantum menor limita cada surto, mas prioridades e empates podem produzir execuções consecutivas do mesmo PID. A recomendação é específica desta carga, desta política e do limiar escolhido; outra carga ou outra meta pode justificar outro quantum.

\clearpage
\section*{6. Política, implementação e reprodução}
\subsection*{Estruturas e decisões de escalonamento}
O BCP armazena PID, nome/programa, PC, estado, prioridade, créditos, X/Y e espera. A tabela possui os BCPs; prontos e bloqueados guardam somente ProcessId. A fila de prontos é ordenada por créditos, e bloqueados seguem ordem de chegada.

Créditos começam na prioridade e diminuem em uma unidade por despacho, sem ficar negativos. A redistribuição usa as prioridades de todos os processos ativos, inclusive bloqueados, somente quando todos têm zero créditos. Quando prontos têm zero e bloqueados ainda têm crédito, aplica-se FIFO aos prontos.

E/S conta como instrução ao começar, bloqueia por dois surtos posteriores e retoma na instrução seguinte. O surto que inicia a E/S não desconta sua própria espera. Sem prontos, a espera avança por ticks lógicos. SAIDA também conta, remove o processo e registra X/Y finais.

\subsection*{Desempate e custo do contexto}
Na carga inicial, empates seguem a ordem dos arquivos. Após quantum, a reinserção ocorre antes dos créditos iguais; no retorno de E/S, depois deles, respeitando as exceções FIFO. Manter o processo que já rodava quando seus créditos empatam evita salvar e carregar outro contexto em uma máquina real e pode preservar dados em cache. Esses custos não são modelados pelo EP. Mesmo ao repetir o PID, o novo surto conta em D conforme a convenção exigida.

\subsection*{Reprodução e alcance}
Execute o binário Escalonador com programas/ no diretório corrente. Para regenerar este relatório, instale a dependência em requirements-relatorio.txt e execute python3 gerar\_relatorio.py de qualquer diretório. O gerador lê os dados relativos à sua própria localização, valida as duas coleções de logs e só então produz relatorio\_ep1.pdf e relatorio\_ep1.tex. A importação e os testes de validação usam apenas a biblioteca padrão de Python.

gerar\_relatorio.py é a fonte comum do relatório; ReportLab desenha os gráficos vetoriais no PDF e o TEX contém os mesmos pontos em desenhos autossuficientes. O enunciado ep1.pdf fundamenta métricas, regras e material de entrega. O relatório avalia os artefatos locais; não comprova a data ou a realização da submissão via eDisciplinas. A entrega exige um ZIP com número USP, código, logs e este PDF, enviado por um único integrante.

\end{document}
